% GENERATED FILE. Edit canonical lessons, then run npm run generate:books.
% canonical-source-hash: fnv1a64:23200c936cdca1a6
% canonical-lessons: AR-C301-istadama, AR-C301-inqalaba, AR-C301-dahiyya, AR-C301-jarih, AR-C301-asafa-aid

\chapter{An accident on the road}
\label{ch:ar-an-accident-on-the-road}
\hlchaptermodality{301}

\begin{chapteropening}
\textbf{By the end of this chapter:} \emph{I can say to collide, to overturn, a victim, a wounded person and to give first aid.}

Complete the last lesson of chapter 301: I can say to collide, to overturn, a victim, a wounded person and to give first aid.
\end{chapteropening}

\section[\texorpdfstring{\ar{اصطدم} (iṣṭadama)}{اصطدم (iṣṭadama)}]{\ar{اصطدم} (iṣṭadama) — to collide, to crash into}

\label{lesson:AR-C301-istadama}



\begin{quote}
Before the new one: say the Arabic for to steal, then the Arabic for to arrest.
\end{quote}

\subsection*{You'll want to know: \ar{اصطدم}}
\textbf{\ar{اصطدم}} — \emph{iṣṭadama} — \textquotedblleft{}to collide, to crash into\textquotedblright{}.

\begin{quote}
\ar{اصطدمت} \ar{السيارة} \ar{بشجرة} \ar{كبيرة}.
\end{quote}

\begin{quote}
\emph{iṣṭadamat as-sayyāra bi-shajara kabīra.}
\end{quote}

\begin{quote}
The car crashed into a big tree.
\end{quote}

Present: \ar{يصطدم}, \emph{yaṣṭadim}.

It takes \ar{بـ} for what is hit: \ar{اصطدم} \ar{بالجدار}.

\begin{cousinweb}
\textbf{\ar{اصطدم}} (\emph{iṣṭadama}) is built on the native Arabic root \emph{ṣ-d-m} (Lane), \textquotedblleft{}to strike, to knock against\textquotedblright{}. It is Form VIII, where the \emph{t} turns to \emph{ṭ} after \emph{ṣ}: \textquotedblleft{}to strike against each other\textquotedblright{}.
\end{cousinweb}

\begin{grammarlens}[title={What you've built}]
One more word for telling what happened.
\end{grammarlens}

\subsection*{Your turn}
\begin{itemize}
\raggedright
  \item \emph{Say it:} \emph{iṣṭadama}
  \item \emph{Say it:} \emph{iṣṭadama}, once more
  \item \emph{Say it:} \emph{iṣṭadamat as-sayyāra bi-shajara kabīra.} — The car crashed into a big tree.
  \item \emph{Recall:} say the Arabic for to steal, then the Arabic for to arrest, then say \emph{iṣṭadama} again
\end{itemize}

\begin{tcolorbox}[breakable,colback=teal!4,colframe=teal!35!black,title={Before you move on}]
What does \ar{اصطدم} mean? (\textquotedblleft{}To collide, to crash into\textquotedblright{}.) Say it once more.
\end{tcolorbox}
\section[\texorpdfstring{\ar{انقلب} (inqalaba)}{انقلب (inqalaba)}]{\ar{انقلب} (inqalaba) — to overturn, to turn over}

\label{lesson:AR-C301-inqalaba}



\begin{quote}
Before the new one: say the Arabic for to arrest, then the Arabic for to collide.
\end{quote}

\subsection*{You'll want to know: \ar{انقلب}}
\textbf{\ar{انقلب}} — \emph{inqalaba} — \textquotedblleft{}to overturn, to turn over\textquotedblright{}.

\begin{quote}
\ar{انقلبت} \ar{الشاحنة} \ar{على} \ar{جانب} \ar{الطريق}.
\end{quote}

\begin{quote}
\emph{inqalabat ash-shāḥina ʿalā jānib aṭ-ṭarīq.}
\end{quote}

\begin{quote}
The lorry overturned at the side of the road.
\end{quote}

Present: \ar{ينقلب}, \emph{yanqalib}.

\begin{cousinweb}
\textbf{\ar{انقلب}} (\emph{inqalaba}) is built on the native Arabic root \emph{q-l-b} (Lane), \textquotedblleft{}to turn a thing over\textquotedblright{}, the root of \emph{qalb}, the heart, which turns. It is Form VII, \textquotedblleft{}to be turned over\textquotedblright{}.
\end{cousinweb}

\begin{grammarlens}[title={What you've built}]
One more word for telling what happened.
\end{grammarlens}

\subsection*{Your turn}
\begin{itemize}
\raggedright
  \item \emph{Say it:} \emph{inqalaba}
  \item \emph{Say it:} \emph{inqalaba}, once more
  \item \emph{Say it:} \emph{inqalabat ash-shāḥina ʿalā jānib aṭ-ṭarīq.} — The lorry overturned at the side of the road.
  \item \emph{Recall:} say the Arabic for to arrest, then the Arabic for to collide, then say \emph{inqalaba} again
\end{itemize}

\begin{tcolorbox}[breakable,colback=teal!4,colframe=teal!35!black,title={Before you move on}]
What does \ar{انقلب} mean? (\textquotedblleft{}To overturn, to turn over\textquotedblright{}.) Say it once more.
\end{tcolorbox}
\section[\texorpdfstring{\ar{ضحية} (ḍaḥiyya)}{ضحية (ḍaḥiyya)}]{\ar{ضحية} (ḍaḥiyya) — a victim}

\label{lesson:AR-C301-dahiyya}



\begin{quote}
Before the new one: say the Arabic for to collide, then the Arabic for to overturn.
\end{quote}

\subsection*{You'll want to know: \ar{ضحية}}
\textbf{\ar{ضحية}} — \emph{ḍaḥiyya} — \textquotedblleft{}a victim\textquotedblright{}.

\begin{quote}
\ar{كان} \ar{الولد} \ar{ضحية} \ar{الحادث}.
\end{quote}

\begin{quote}
\emph{kāna al-walad ḍaḥiyyat al-ḥādith.}
\end{quote}

\begin{quote}
The boy was the victim of the accident.
\end{quote}

Plural: \ar{ضحايا}, \emph{ḍaḥāyā}.

\begin{cousinweb}
\textbf{\ar{ضحية}} (\emph{ḍaḥiyya}) is built on the native Arabic root \emph{ḍ-ḥ-w} (Lane), the root of the taught \emph{ḍuḥā}, the bright late morning. Lane gives \emph{ḍaḥiyya} first as an animal offered at that hour; today it is anyone who suffers harm.
\end{cousinweb}

\begin{grammarlens}[title={What you've built}]
One more word for telling what happened.
\end{grammarlens}

\subsection*{Your turn}
\begin{itemize}
\raggedright
  \item \emph{Say it:} \emph{ḍaḥiyya}
  \item \emph{Say it:} \emph{ḍaḥiyya}, once more
  \item \emph{Say it:} \emph{kāna al-walad ḍaḥiyyat al-ḥādith.} — The boy was the victim of the accident.
  \item \emph{Recall:} say the Arabic for to collide, then the Arabic for to overturn, then say \emph{ḍaḥiyya} again
\end{itemize}

\begin{tcolorbox}[breakable,colback=teal!4,colframe=teal!35!black,title={Before you move on}]
What does \ar{ضحية} mean? (\textquotedblleft{}A victim\textquotedblright{}.) Say it once more.
\end{tcolorbox}
\section[\texorpdfstring{\ar{جريح} (jarīḥ)}{جريح (jarīḥ)}]{\ar{جريح} (jarīḥ) — a wounded person; wounded}

\label{lesson:AR-C301-jarih}



\begin{quote}
Before the new one: say the Arabic for to overturn, then the Arabic for a victim.
\end{quote}

\subsection*{You'll want to know: \ar{جريح}}
\textbf{\ar{جريح}} — \emph{jarīḥ} — \textquotedblleft{}a wounded person; wounded\textquotedblright{}.

\begin{quote}
\ar{نقلت} \ar{سيارة} \ar{الإسعاف} \ar{الجريح} \ar{إلى} \ar{المستشفى}.
\end{quote}

\begin{quote}
\emph{naqalat sayyārat al-isʿāf al-jarīḥ ilā al-mustashfā.}
\end{quote}

\begin{quote}
The ambulance took the wounded man to hospital.
\end{quote}

Plural: \ar{جرحى}, \emph{jarḥā}.

\begin{cousinweb}
\textbf{\ar{جريح}} (\emph{jarīḥ}) is built on the native Arabic root \emph{j-r-ḥ} (Lane), \textquotedblleft{}to wound\textquotedblright{}, the root of the taught \emph{jurḥ}, a wound. The \emph{faʿīl} pattern gives the one who bears it.
\end{cousinweb}

\begin{grammarlens}[title={What you've built}]
One more word for telling what happened.
\end{grammarlens}

\subsection*{Your turn}
\begin{itemize}
\raggedright
  \item \emph{Say it:} \emph{jarīḥ}
  \item \emph{Say it:} \emph{jarīḥ}, once more
  \item \emph{Say it:} \emph{naqalat sayyārat al-isʿāf al-jarīḥ ilā al-mustashfā.} — The ambulance took the wounded man to hospital.
  \item \emph{Recall:} say the Arabic for to overturn, then the Arabic for a victim, then say \emph{jarīḥ} again
\end{itemize}

\begin{tcolorbox}[breakable,colback=teal!4,colframe=teal!35!black,title={Before you move on}]
What does \ar{جريح} mean? (\textquotedblleft{}A wounded person; wounded\textquotedblright{}.) Say it once more.
\end{tcolorbox}
\section[\texorpdfstring{\ar{أسعف} (asʿafa)}{أسعف (asʿafa)}]{\ar{أسعف} (asʿafa) — to give first aid}

\label{lesson:AR-C301-asafa-aid}



\begin{quote}
Before the new one: say the Arabic for a victim, then the Arabic for a wounded person.
\end{quote}

\subsection*{You'll want to know: \ar{أسعف}}
\textbf{\ar{أسعف}} — \emph{asʿafa} — \textquotedblleft{}to give first aid\textquotedblright{}.

\begin{quote}
\ar{أسعف} \ar{الطبيب} \ar{الولد} \ar{في} \ar{الشارع}.
\end{quote}

\begin{quote}
\emph{asʿafa aṭ-ṭabīb al-walad fī ash-shāriʿ.}
\end{quote}

\begin{quote}
The doctor gave the boy first aid in the street.
\end{quote}

Present: \ar{يسعف}, \emph{yusʿif}.

The taught \ar{إسعاف}, first aid, is its verbal noun.

\begin{cousinweb}
\textbf{\ar{أسعف}} (\emph{asʿafa}) is built on the native Arabic root \emph{s-ʿ-f} (Lane), \textquotedblleft{}to help, to give someone what he needs\textquotedblright{}. It is Form IV; \emph{isʿāf} is its verbal noun.
\end{cousinweb}

\begin{grammarlens}[title={What you've built}]
One more word for telling what happened.
\end{grammarlens}

\subsection*{Your turn}
\begin{itemize}
\raggedright
  \item \emph{Say it:} \emph{asʿafa}
  \item \emph{Say it:} \emph{asʿafa}, once more
  \item \emph{Say it:} \emph{asʿafa aṭ-ṭabīb al-walad fī ash-shāriʿ.} — The doctor gave the boy first aid in the street.
  \item \emph{Recall:} say the Arabic for a victim, then the Arabic for a wounded person, then say \emph{asʿafa} again
\end{itemize}

\begin{tcolorbox}[breakable,colback=teal!4,colframe=teal!35!black,title={Before you move on}]
What does \ar{أسعف} mean? (\textquotedblleft{}To give first aid\textquotedblright{}.) Say it once more.
\end{tcolorbox}
